\chapter{TỔNG QUAN ĐỀ TÀI}
\label{ch:01}

\section{Bối cảnh bài toán quản lý kho tại các doanh nghiệp vừa và nhỏ}
Trong chuỗi cung ứng hiện đại, kho bãi giữ vai trò mạch máu lưu chuyển dòng vật chất và vốn của doanh nghiệp. Qua quá trình khảo sát thực tế tại các đơn vị phân phối linh kiện máy tính và thiết bị mạng, các doanh nghiệp vừa và nhỏ (SMEs) thường gặp phải các hạn chế nghiêm trọng do sử dụng sổ tay hoặc Excel phân tán: độ trễ thông tin từ 1 đến 3 ngày, nguy cơ sai sót cao và tình trạng chôn vốn lưu động.

\section{Mục tiêu, đối tượng và phạm vi nghiên cứu}
Đề tài hướng tới xây dựng hệ thống Quản lý Kho thông minh có tích hợp Trí tuệ Nhân tạo, tự động hóa xuất nhập tồn, phân tích vận tốc tiêu thụ và cung cấp trợ lý hỏi đáp tự nhiên. Đối tượng áp dụng là các nhà kho phân phối quy mô vừa với phân quyền 3 vai trò: Quản trị viên, Thủ kho và Kế toán.

\section{Phương pháp nghiên cứu và cách tiếp cận ứng dụng AI}
Ứng dụng nguyên tắc Chống ảo giác dữ liệu (Anti-Hallucination) với Ground-Truth Context Injection và xây dựng kiến trúc dự phòng kép (Dual-layer Resilience) đảm bảo hệ thống luôn hoạt động liên tục.
